\label{chap:methodology}

\section{System Architecture Overview}
The Soft L-Topological architecture works for both causal text generation and bidirectional language tasks. Rather than computing dot products in Euclidean space and normalizing them with softmax, we use \textbf{bounded soft set representations} and \textbf{nilpotent lattice routing}.

The forward pass processes inputs in four main steps:
\begin{enumerate}[itemsep=1.5pt, topsep=2pt, parsep=0pt]
    \item Tokens $w_t \in \mathcal{V}$ are projected directly into the unit cube $[0, 1]^D$.
    \item Learnable positional vectors in $[0, 1]^D$ are combined with token embeddings using convex blending.
    \item A stack of $L$ \textbf{Lattice Blocks} transforms the states using multi-head lattice routing, residual connections, and position-wise feed-forward networks (FFN).
    \item A linear output head converts hidden representations into vocabulary logits.
\end{enumerate}

\section{Bounded Continuous Lattice Embeddings}
Because fuzzy operators are defined on the interval $[0, 1]$, every token representation must stay strictly inside the unit cube $[0, 1]^D$ without drifting into unbounded values.

\subsection{Convex Token and Positional Blending}
We pass raw token weights $\widetilde{E} \in \mathbb{R}^{|\mathcal{V}| \times D}$ and position weights $\widetilde{P} \in \mathbb{R}^{S \times D}$ through a sigmoid function: $\mathbf{e}_{\text{tok}}(w) = \sigma(\widetilde{E}_w) \in (0, 1)^D$ and $\mathbf{e}_{\text{pos}}(t) = \sigma(\widetilde{P}_t) \in (0, 1)^D$. If we simply added them together, the values would exceed 1.0. Instead, we combine them through \textbf{convex blending} with $\lambda = 0.5$:
\begin{equation}
\label{eq:convex_blending}
\mathbf{x}_t^{(0)} = \lambda \mathbf{e}_{\text{tok}}(w_t) + (1 - \lambda) \mathbf{e}_{\text{pos}}(t) \in [0, 1]^D
\end{equation}
Because the unit hypercube $[0, 1]^D$ is convex, $\mathbf{x}_t^{(0)} \in [0, 1]^D$ holds strictly for all sequence positions.

\subsection{Beta Distribution Initialization}
Standard Gaussian initialization centered at zero causes the sigmoid function to map almost every weight close to $0.5$. If all attributes start at $0.5$, tokens share nearly identical overlap at the beginning of training. To give the model a diverse spread of initial attribute memberships, we initialize embedding weights using an asymmetric \textbf{Beta distribution}:
\begin{equation}
\widetilde{E}_{vd} \sim \text{logit}(\text{Beta}(\alpha, \beta)), \quad \alpha = 1.0, \, \beta = 3.0
\end{equation}
This concentrates initial memberships between $0.1$ and $0.4$, setting most attribute pairs near or below the nilpotent threshold at initialization.

\section{The Continuous Lattice Sequence Router}
For input states $X \in \mathbb{R}^{S \times D}$, the router splits representations into $H$ parallel heads of dimension $d_k = D / H$: Queries $Q^{(h)} = \sigma(X W_Q^{(h)}) \in [0, 1]^{S \times d_k}$, Keys $K^{(h)} = \sigma(X W_K^{(h)}) \in [0, 1]^{S \times d_k}$, and Values $V^{(h)} = X W_V^{(h)} \in \mathbb{R}^{S \times d_v}$. Queries and Keys serve as bounded truth values in $[0, 1]^{d_k}$ for routing, while Values remain in unconstrained $\mathbb{R}^{d_v}$ so the model retains full capacity for semantic features.

\subsection{Thresholded \L{}ukasiewicz Conjunction}
For each head $h$, the multi-attribute semantic conjunction between query token $i$ and key token $j$ is computed using the thresholded \L{}ukasiewicz $t$-norm with activation threshold $\tau \ge 0$:
\begin{equation}
\label{eq:thresholded_lukasiewicz}
\mathcal{C}_{ij}^{(h)} = \sum_{d=1}^{d_k} \max(0, Q_{id}^{(h)} + K_{jd}^{(h)} - 1 - \tau)
\end{equation}
If $Q_{id}^{(h)} + K_{jd}^{(h)} \le 1 + \tau$ across all coordinates, the total overlap $\mathcal{C}_{ij}^{(h)}$ evaluates to exactly \textbf{0.0000}.

\subsection{Continuous Jaccard Relative Containment}
To normalize routing weights between 0 and 1 without risking division by zero, we compute the \textbf{continuous Jaccard containment coefficient} by dividing conjunction by the continuous fuzzy union ($t$-conorm):
\begin{equation}
\label{eq:jaccard_routing}
R_{ij}^{(h)} = \frac{\sum_{d=1}^{d_k} \max(0, Q_{id}^{(h)} + K_{jd}^{(h)} - 1 - \tau)}{\sum_{d=1}^{d_k} \left(Q_{id}^{(h)} + K_{jd}^{(h)} - Q_{id}^{(h)} K_{jd}^{(h)}\right) + \epsilon}
\end{equation}
where $\epsilon = 10^{-7}$. This formulation ensures that $0.0000 \le R_{ij}^{(h)} \le 1.0000$ always, and $R_{ij}^{(h)} \equiv 0.0000$ whenever the numerator is zero.

\subsection{Causal Masking and Value Aggregation}
For autoregressive generation, future connections are set directly to zero ($R_{ij}^{(h)} = 0.0000$ for $j > i$). Output representations are formed via value aggregation followed by multi-head projection and a standard two-layer GELU feed-forward network:
\begin{equation}
\text{Head}^{(h)}_i = \sum_{j=1}^i R_{ij}^{(h)} V_j^{(h)}, \quad X^{(l+1)} = X^{(l)\prime} + \text{FFN}(\text{LayerNorm}(X^{(l)\prime}))
\end{equation}

\section{Continuous Lattice KV-Caching}
Generating tokens autoregressively without caching recomputes the entire history at every step, requiring $O(S^2)$ operations. With our Key-Value cache, we store previous keys $K_{\text{past}} \in [0, 1]^{t \times d_k}$ and values $V_{\text{past}} \in \mathbb{R}^{t \times d_v}$. At step $t$, the new query $q_t$ only needs to evaluate Equation \ref{eq:jaccard_routing} against the cached keys, lowering decoding complexity to $O(T)$.

If a token has almost zero overlap with every past token as well as itself ($\sum_{j \le t} R_{tj} < 10^{-6}$), the router applies an \textbf{exact causal diagonal fallback}: it sets $R_{tt} = 1.0$, allowing the token to attend to itself. This prevents numerical instability when denominator sums are tiny. In our automated tests, cached output matched uncached output down to machine epsilon: $\Delta = \max | X_{\text{cached}} - X_{\text{uncached}} | = 5.06 \times 10^{-7}$.
